\chapter{The Squeakymail Example}

This chapter will present a simple yet complete Common Lisp\index{Common Lisp} application. The application
is a web-based email client, called ``Squeakymail,'' which can be used for reading, composing,
sending, and managing email through a web browser. The complete code for this example
can be downloaded from

\begin{center}
\texttt{http://www.genworks.com/downloads}
\end{center}

In this chapter, we will give an overview of the application with excerpts from the source
code, to help tie together the concepts covered in earlier chapters. This example is intended
to give a flavor of working with CL on a real application.

\section{Requirements and Design Overview}

The basic requirements for this application are the following:

\begin{itemize}
\item Must be accessible from any web browser
\item Must handle POP3 mailboxes
\item Should allow the viewing of mail messages
\item Should allow the composition and sending of mail
\item Should allow the deletion and management of messages
\end{itemize}

Fortunately, Franz Inc. provides a free open-source library for handling POP3 email and
for composing and sending mail. This library can also be downloaded from the above URL.

The webserver will use AllegroServe, as introduced in Chapter 4. The overall design
breaks down as follows:

\begin{enumerate}
\item The AllegroServe session will store a list of messages in a global parameter
\item A background thread will periodically poll the POP3 server and update the list with
any new mail
\item The web interface will provide pages to view the list of messages, to read the content
of individual messages, and to compose and send new mail
\item Sufficient security will be in place to prevent multiple users of the application from
reading each other's email
\end{enumerate}

\section{Main Application File}

The main file for the application looks like this:

\begin{verbatim}
(in-package :user)

(defparameter *squeaky-port* 9876)

(defparameter *imap-server* "localhost")
(defparameter *pop-server* "localhost")
(defparameter *smtp-server* "localhost")

(defparameter *username* (sys:user-name))

(defparameter *poll-interval* 30)

(defun start-squeaky (&key (port *squeaky-port*))
  (start-server port)
  (start-background-polling))

(defun stop-squeaky ()
  (stop-background-polling)
  (stop-server))
\end{verbatim}

The function \texttt{start-squeaky} starts the webserver and kicks off a background thread to
poll the mailbox server for new mail. The \texttt{stop-squeaky} function stops this background
thread and shuts down the webserver.

\section{Working With Global State}

Because the web application runs in a persistent CL session, the entire program state can
be maintained in memory for fast and convenient access. No need to maintain state in a
database or in some transient or persistent file system.

For this application, we maintain the state in three main global variables:

\begin{verbatim}
(defparameter *lock* (mp:make-process-lock))

(defparameter *messages* nil)

(defparameter *pollingp* nil)
\end{verbatim}

The \texttt{*lock*} variable is a process lock, used for controlling access to the \texttt{*messages*} list
in a multi-threaded environment. The \texttt{*messages*} variable contains a list of mail-message
objects. The \texttt{*pollingp*} variable serves as a Boolean flag indicating whether the background
polling process is currently running.

\section{Background Polling}

The \texttt{start-background-polling} and \texttt{stop-background-polling} functions spawn and
halt a background thread which periodically polls for new mail:

\begin{verbatim}
(defun start-background-polling ()
  (setq *pollingp* t)
  (mp:process-run-function "Polling Process"
    #'polling-function))

(defun stop-background-polling ()
  (setq *pollingp* nil))

(defun polling-function ()
  (do () ((not *pollingp*))
    (update-messages-from-server)
    (sleep *poll-interval*)))
\end{verbatim}

The call to \texttt{mp:process-run-function} spawns a new operating system thread. This function takes a name for the thread, and a function which the thread should invoke. The
polling thread simply calls the \texttt{update-messages-from-server} function in a loop, waiting \texttt{*poll-interval*} seconds between each iteration.

The \texttt{update-messages-from-server} function is the function which actually contacts
the POP3 server:

\begin{verbatim}
(defun update-messages-from-server ()
  (mp:with-process-lock (*lock*)
    (let ((count (ignore-errors (pop:pop-message-count *pop-server*
                                                       *username*))))
      (when (and count (numberp count))
        (dotimes (n count)
          (unless (find n *messages* :key #'message-number)
            (let ((message (retrieve-message (1+ n))))
              (push message *messages*))))))))
\end{verbatim}

The \texttt{mp:with-process-lock} macro ensures that only one thread at a time can be accessing
and modifying the \texttt{*messages*} list.

The \texttt{pop:pop-message-count} function returns the number of messages currently waiting on the server for the given username. The username in this example defaults to the
Unix login name for the user running the CL session. You will have to ensure that you can
authenticate yourself to your POP3 server with this username.

If there are messages on the server which are not yet in the \texttt{*messages*} list, they are
retrieved and added to the list. The details of the \texttt{retrieve-message} function, which
wraps the calls to retrieve the actual message, are not shown.

\section{Serving the Web Pages}

The application provides three main web pages:

\begin{enumerate}
\item Index --- shows a list of the available messages, with links to read them
\item Read --- shows the content of an individual message
\item Compose --- allows the user to compose and send a new message
\end{enumerate}

Here is the function which publishes the main index page:

\begin{verbatim}
(publish :path "/squeaky"
         :content-type "text/html"
         :function
         #'(lambda (req ent)
             (with-http-response (req ent)
               (with-http-body (req ent)
                 (html
                  (:html
                   (:head (:title "Squeakymail"))
                   (:body
                    (:h1 "Your Mail")
                    ((:form :method "post" :action "/squeaky/compose")
                     ((:input :type "submit" :value "Compose New")))
                    (:hr)
                    (if *messages*
                        (html
                         ((:table :border 1)
                          (:tr (:th "Number") (:th "From")
                               (:th "Subject") (:th "Date"))
                          (dolist (message (reverse *messages*))
                            (html
                             (:tr
                              (:td ((:a :href
                                     (format nil "/squeaky/read?number=~a"
                                             (message-number message)))
                                    (:princ (message-number message))))
                              (:td (:princ (message-from message)))
                              (:td (:princ (message-subject message)))
                              (:td (:princ (message-date message))))))))
                        (html (:i "No messages")))
                    (:hr))))))))
\end{verbatim}

This function uses the \texttt{html} macro from AllegroServe to generate the HTML content
dynamically. The page includes a button to compose new mail, and a table showing all the
messages currently in the \texttt{*messages*} list. Each message number is a hyperlink to read that
particular message.

The ``read'' page is published in a similar manner:

\begin{verbatim}
(publish :path "/squeaky/read"
         :content-type "text/html"
         :function
         #'(lambda (req ent)
             (let* ((number-string (request-query-value "number" req))
                    (number (when number-string
                              (ignore-errors (parse-integer number-string))))
                    (message (when number
                               (find number *messages*
                                     :key #'message-number))))
               (with-http-response (req ent)
                 (with-http-body (req ent)
                   (if message
                       (html
                        (:html
                         (:head (:title "Read Mail"))
                         (:body
                          (:h2 "Message " (:princ number))
                          (:b "From: ") (:princ (message-from message)) (:br)
                          (:b "Subject: ") (:princ (message-subject message))
                          (:br)
                          (:b "Date: ") (:princ (message-date message)) (:br)
                          (:hr)
                          (:pre (:princ (message-body message)))
                          (:hr)
                          ((:form :method "post" :action "/squeaky/delete")
                           ((:input :type "hidden" :name "number"
                                    :value (format nil "~a" number)))
                           ((:input :type "submit" :value "Delete")))
                          ((:a :href "/squeaky") "Back to Index"))))
                       (html
                        (:html
                         (:head (:title "Error"))
                         (:body
                          (:h2 "Message not found")
                          ((:a :href "/squeaky") "Back to Index"))))))))))
\end{verbatim}

This page retrieves the message number from the query string, finds the corresponding
message in the \texttt{*messages*} list, and displays its content. It also provides a button to
delete the message and a link back to the main index.

\section{Composing and Sending Mail}

The compose page is slightly more complex, as it needs to handle both displaying a blank
form and processing the submitted form data:

\begin{verbatim}
(publish :path "/squeaky/compose"
         :content-type "text/html"
         :function
         #'(lambda (req ent)
             (let ((to (request-query-value "to" req))
                   (subject (request-query-value "subject" req))
                   (body (request-query-value "body" req)))
               (with-http-response (req ent)
                 (with-http-body (req ent)
                   (if (and to subject body)
                       (progn
                         (send-smtp-mail *smtp-server* *username*
                                        to subject body)
                         (html
                          (:html
                           (:head (:title "Message Sent"))
                           (:body
                            (:h2 "Your message has been sent")
                            ((:a :href "/squeaky") "Back to Index")))))
                       (html
                        (:html
                         (:head (:title "Compose Mail"))
                         (:body
                          (:h2 "Compose New Message")
                          ((:form :method "post" :action "/squeaky/compose")
                           (:b "To: ")
                           ((:input :type "text" :name "to" :size 40)) (:br)
                           (:b "Subject: ")
                           ((:input :type "text" :name "subject" :size 40))
                           (:br) (:br)
                           ((:textarea :name "body" :rows 20 :cols 60))
                           (:br) (:br)
                           ((:input :type "submit" :value "Send"))
                           "&nbsp;&nbsp;"
                           ((:input :type "button" :value "Cancel"
                                    :onclick "window.location='/squeaky'")))
                          (:hr)))))))))))
\end{verbatim}

If the form has been submitted (i.e., if the ``to,'' ``subject,'' and ``body'' parameters are present),
the \texttt{send-smtp-mail} function is called to send the message. Otherwise, a blank form is
displayed.

\section{Deleting Messages}

The delete function removes a message from the local list and also deletes it from the POP3
server:

\begin{verbatim}
(publish :path "/squeaky/delete"
         :content-type "text/html"
         :function
         #'(lambda (req ent)
             (let* ((number-string (request-query-value "number" req))
                    (number (when number-string
                              (ignore-errors (parse-integer number-string)))))
               (when number
                 (mp:with-process-lock (*lock*)
                   (setq *messages* (remove number *messages*
                                           :key #'message-number))
                   (ignore-errors
                     (pop:delete-letter *pop-server* *username* number))))
               (with-http-response (req ent)
                 (with-http-body (req ent)
                   (html
                    (:html
                     (:head (:title "Message Deleted"))
                     (:body
                      (:h2 "Message has been deleted")
                      ((:a :href "/squeaky") "Back to Index")))))))))
\end{verbatim}

Note the use of \texttt{mp:with-process-lock} to ensure thread-safe modification of the \texttt{*messages*}
list.

\section{Starting the Webserver}

The \texttt{start-server} and \texttt{stop-server} functions are simple wrappers around the AllegroServe functions:

\begin{verbatim}
(defun start-server (port)
  (net.aserve:start :port port))

(defun stop-server ()
  (net.aserve:shutdown))
\end{verbatim}

\section{Message Data Structure}

The application uses a simple structure to represent mail messages:

\begin{verbatim}
(defstruct message
  number
  from
  subject
  date
  body)
\end{verbatim}

The \texttt{defstruct}\index{defstruct@\texttt{defstruct}} macro automatically creates constructor, accessor, and type-predicate functions for this structure.

\section{Testing the Application}

To test Squeakymail, you will need:

\begin{enumerate}
\item A running POP3 mail server (or access to one)
\item A running SMTP mail server (or access to one)
\item The Post Office library from Franz Inc.
\item AllegroServe
\end{enumerate}

Once these components are in place, you can start the application by loading all the source
files and calling:

\begin{verbatim}
(start-squeaky)
\end{verbatim}

Then point your web browser to:

\begin{verbatim}
http://localhost:9876/squeaky
\end{verbatim}

You should see the main index page with your mail messages listed. From there you can
read, compose, and delete messages.

\section{Security Considerations}

This example application is intended for demonstration purposes only and should not be used
in a production environment without significant security enhancements. In particular:

\begin{itemize}
\item The application does not implement user authentication
\item Passwords are not encrypted
\item There is no session management to prevent unauthorized access
\item The application assumes a single user
\end{itemize}

For a production application, you would need to implement proper authentication, session
management, and access controls.

\section{Enhancements}

Some possible enhancements to this application might include:

\begin{itemize}
\item Support for IMAP in addition to POP3
\item Ability to reply to messages
\item Support for attachments
\item Better formatting of HTML email
\item Address book functionality
\item Folders and message filtering
\item Search functionality
\item Support for multiple users with authentication
\end{itemize}

\section{Conclusion}

This chapter has demonstrated a complete, albeit simple, web application written in Common
Lisp. The example shows how the various concepts covered in earlier chapters come together
in a real application:

\begin{itemize}
\item Use of global parameters to maintain state
\item Multi-threaded programming with process locks
\item Web serving with AllegroServe
\item Dynamic HTML generation
\item External library integration (POP3/SMTP)
\item Data structures and list manipulation
\end{itemize}

The complete source code for Squeakymail, along with installation instructions, can be
downloaded from the URL given at the beginning of this chapter. We encourage you to
download, install, and experiment with this application as a learning tool for Common Lisp
development.
